\exercise{10.9}{A decreasing recurrence with a squared term}

\goal For
\[
 s_1=1,\qquad s_{n+1}=\frac{n}{n+1}s_n^2\quad(n\in\N),
\]
compute $s_2,s_3,s_4$, prove that a finite limit exists, and prove
that its value is $0$.

\step{Part (a): substitute the indices one at a time.}
For $n=1$, the recurrence gives
\[
 s_2=\frac1{1+1}s_1^2=\frac12\cdot1^2=\frac12.
\]
For $n=2$,
\[
 s_3=\frac2{2+1}s_2^2
     =\frac23\left(\frac12\right)^2
     =\frac23\cdot\frac14=\frac16.
\]
For $n=3$,
\[
 s_4=\frac3{3+1}s_3^2
     =\frac34\left(\frac16\right)^2
     =\frac34\cdot\frac1{36}
     =\frac3{144}=\frac1{48}.
\]

\answer{\(\boxed{s_2=\tfrac12,\quad s_3=\tfrac16,\quad
s_4=\tfrac1{48}}\).}

\step{Part (b), Step 1: prove positivity and an upper bound by induction.}
We claim $0<s_n\le1$ for every $n\in\N$.

The initial case is $s_1=1$, so $0<s_1\le1$ holds.
For the induction step, assume $0<s_n\le1$.
Since $n\ge1$, the coefficient satisfies
\[
 0<\frac n{n+1}<1.
\]
Also $0<s_n^2\le1$.
Multiplying these positive quantities gives
\[
 0<s_{n+1}=\frac n{n+1}s_n^2<1.
\]
This proves the claim by induction.

\step{Part (b), Step 2: compare consecutive terms.}
The inequality $s_n\le1$, multiplied by the positive number $s_n$,
gives $s_n^2\le s_n$. The coefficient $n/(n+1)$ is strictly below $1$,
and $s_n^2>0$, so
\[
 s_{n+1}=\frac n{n+1}s_n^2<s_n^2\le s_n.
\]
Hence the sequence is strictly decreasing. It is bounded below by $0$
and above by $1$.

\step{Part (b), Step 3: apply monotone convergence before manipulating a limit.}
Theorem~10.2 for bounded monotone sequences now applies.
There is a real number $L$ with $s_n\to L$.
The inequalities $s_n>0$ imply $L\ge0$ by Exercise~8.9(a).
Since the sequence is decreasing and $s_2=1/2$,
\[
 0<s_n\le\frac12\qquad(n\ge2).
\]
Exercise~8.9 also passes this eventual upper bound to the limit.
Therefore
\[
 0\le L\le\frac12.
\]

\answer{A finite limit exists, and it belongs to $[0,1/2]$.
Its existence has been proved independently of any equation for its value.}

\step{Part (c), Step 1: justify the shifted limit and the coefficient limit.}
If $s_n\to L$, then $s_{n+1}\to L$ as well.
Indeed, for any $\varepsilon>0$ choose a threshold $N$ with
$m>N\Rightarrow\abs{s_m-L}<\varepsilon$.
For $n>N$, we have $n+1>N$, so the same estimate holds with $m=n+1$.

Next put $c_n=n/(n+1)$. Its error from $1$ is
\[
 \abs{c_n-1}
 =\abs{\frac{n-(n+1)}{n+1}}
 =\frac1{n+1}.
\]
For any $\varepsilon>0$, if $n>1/\varepsilon$, then
$n+1>1/\varepsilon$, so $1/(n+1)<\varepsilon$.
Thus $c_n\to1$.

\step{Part (c), Step 2: apply the product limit law to the established limits.}
Because $s_n\to L$ and $c_n\to1$, the product laws from Section~9 give
\[
 \lim_{n\to\infty}c_ns_n^2=1\cdot L\cdot L=L^2.
\]
The recurrence says that this is the same sequence as $s_{n+1}$,
whose limit is $L$. Uniqueness of a finite limit gives
\[
 L=L^2.
\]
Rearranging and factoring gives
\[
 L^2-L=0,\qquad L(L-1)=0.
\]
The zero-product property therefore gives $L=0$ or $L=1$.

\step{Part (c), Step 3: use the tail bound to remove the extraneous candidate.}
Part (b) proved $L\le1/2$, so $L=1$ is impossible.
The only remaining candidate is $L=0$.

\answer{\(\boxed{\displaystyle\lim_{n\to\infty}s_n=0}\).}

\step{An additional direct error bound.}
The same recurrence supplies a quantitative check. Since $s_n^2\le s_n$,
\[
 s_{n+1}\le\frac n{n+1}s_n.
\]
Starting with $s_1=1=1/1$, induction gives $s_n\le1/n$:
if $s_n\le1/n$, then
\[
 s_{n+1}\le\frac n{n+1}\cdot\frac1n=\frac1{n+1}.
\]
Consequently
\[
 \abs{s_n-0}=s_n\le\frac1n<\varepsilon
 \qquad\left(n>\frac1\varepsilon\right).
\]
This also verifies the zero limit directly from Definition~7.1.

\exercise{10.10}{A linear recurrence approaching one half}

\goal For
\[
 s_1=1,\qquad s_{n+1}=\frac{s_n+1}{3}\quad(n\in\N),
\]
compute $s_2,s_3,s_4$, prove $s_n>1/2$ by induction, prove that the
sequence is decreasing, and identify its limit.

\step{Part (a): calculate the first three new terms.}
Using the recurrence successively gives
\[
 s_2=\frac{1+1}{3}=\frac23,
\]
\[
 s_3=\frac{\frac23+1}{3}
     =\frac{\frac23+\frac33}{3}
     =\frac{5/3}{3}=\frac59,
\]
and
\[
 s_4=\frac{\frac59+1}{3}
     =\frac{\frac59+\frac99}{3}
     =\frac{14/9}{3}=\frac{14}{27}.
\]

\answer{\(\boxed{s_2=\tfrac23,\quad s_3=\tfrac59,\quad
s_4=\tfrac{14}{27}}\).}

\step{Part (b), Step 1: verify the initial case of the proposed lower bound.}
The proposition for induction is $P_n$: $s_n>1/2$.
For $n=1$, it says $1>1/2$, which is true.

\step{Part (b), Step 2: prove the induction implication.}
Assume $s_n>1/2$.
Adding $1$ to both sides gives $s_n+1>3/2$.
Dividing by the positive number $3$ preserves the strict inequality:
\[
 s_{n+1}=\frac{s_n+1}{3}>\frac{3/2}{3}=\frac12.
\]
Thus $P_n$ implies $P_{n+1}$. Induction proves the bound for every
positive integer $n$.

\answer{\(\boxed{s_n>\tfrac12\text{ for every }n\in\N}\).}

\step{Part (c): subtract consecutive terms and determine the sign.}
Use a common denominator:
\[
 s_{n+1}-s_n
 =\frac{s_n+1}{3}-s_n
 =\frac{s_n+1-3s_n}{3}
 =\frac{1-2s_n}{3}.
\]
Part (b) gives $2s_n>1$, so the numerator is negative.
The denominator $3$ is positive. Hence $s_{n+1}-s_n<0$,
which means $s_{n+1}<s_n$ for every $n$.

\answer{The sequence is strictly decreasing.}

\step{Part (d), Step 1: establish convergence from bounds and monotonicity.}
The terms satisfy $1/2<s_n\le s_1=1$, and part (c) proves that
the sequence is decreasing. Theorem~10.2 therefore gives a finite
limit $L$. Exercise~8.9 implies
\[
 \frac12\le L\le1.
\]

\step{Part (d), Step 2: solve the limit equation after convergence is known.}
The shifted sequence $s_{n+1}$ has the same limit $L$, by the
threshold argument given in Exercise~10.9(c).
The addition and scalar-multiplication laws from Section~9 give
\[
 \lim_{n\to\infty}\frac{s_n+1}{3}=\frac{L+1}{3}.
\]
The recurrence identifies these two sequences, so their limits agree:
\[
 L=\frac{L+1}{3}.
\]
Multiplying by $3$ and collecting terms gives
\[
 3L=L+1,\qquad 2L=1,\qquad L=\frac12.
\]

\answer{\(\boxed{\displaystyle\lim_{n\to\infty}s_n=\tfrac12}\).
Every term is strictly greater than $1/2$, but the limit can equal
that lower bound.}

\exercise{10.11}{An infinite-product recurrence}

\goal For
\[
 t_1=1,\qquad
 t_{n+1}=\left(1-\frac1{4n^2}\right)t_n\quad(n\in\N),
\]
(a) prove that a finite limit exists and (b) investigate its value.
Part (b) asks what the limit appears to be. We will distinguish
a numerical guess from a proof; an optional elementary-calculus
argument will then verify the exact value.

\step{Part (a), Step 1: bound the coefficient and prove positivity.}
For every $n\ge1$, we have $4n^2\ge4$, so
\[
 0<\frac1{4n^2}\le\frac14.
\]
Subtracting from $1$ gives
\[
 \frac34\le1-\frac1{4n^2}<1.
\]
In particular, every recurrence coefficient is strictly positive.
The initial value is $t_1=1>0$. If $t_n>0$, the recurrence is a
product of two positive numbers and gives $t_{n+1}>0$.
Induction shows $t_n>0$ for all $n$.

\step{Part (a), Step 2: prove monotonicity and apply the theorem.}
Since $1-1/(4n^2)<1$ and $t_n>0$,
\[
 t_{n+1}=\left(1-\frac1{4n^2}\right)t_n<t_n.
\]
Thus $(t_n)$ is strictly decreasing, bounded below by $0$ and above
by its first term $1$. Theorem~10.2 gives a finite real limit $L$.
The bound on the terms implies $0\le L\le1$ by Exercise~8.9.

\answer{The limit of $(t_n)$ exists as a real number.
This conclusion uses only positivity, monotonicity, and boundedness.}

\step{Part (b), Step 1: compute initial values and a numerical table.}
The first terms are
\[
 t_2=\left(1-\frac14\right)1=\frac34,
\]
\[
 t_3=\left(1-\frac1{16}\right)\frac34
     =\frac{15}{16}\cdot\frac34=\frac{45}{64},
\]
and
\[
 t_4=\left(1-\frac1{36}\right)\frac{45}{64}
     =\frac{35}{36}\cdot\frac{45}{64}
     =\frac{1575}{2304}=\frac{175}{256}.
\]
Continuing the recurrence produces the following rounded values:
\[
 \begin{array}{c|c}
 n&t_n\\ \hline
 1&1.000000000000\\
 2&0.750000000000\\
 3&0.703125000000\\
 4&0.683593750000\\
 10&0.653587392299\\
 100&0.638221323951\\
 1000&0.636779026837
 \end{array}
\]
The numbers suggest a limit near $0.63662$.
The familiar constant $2/\pi$ has the decimal value
$0.636619772368\ldots$, which suggests the guess
\[
 L=\frac2\pi.
\]
The numerical agreement motivates this guess; finitely many numerical
calculations by themselves do not prove the exact identification.

\step{Part (b), Step 2: explain why a direct limit equation gives no value.}
The coefficient $1-1/(4n^2)$ tends to $1$.
Taking limits in the recurrence, after convergence has been proved,
only yields
\[
 L=1\cdot L.
\]
Every real $L$ satisfies this identity, so it does not determine the
limit. Further information is needed to justify the guess $2/\pi$.

\step{Optional verification of the guess: the elementary-calculus input.}
The remaining argument uses integration by parts and the order
property of definite integrals. Those are additional calculus tools;
they were not needed for the existence proof in part (a).
We derive the required product formula here rather than assume
Wallis's product.

For an integer $m\ge0$, define
\[
 I_m=\int_0^{\pi/2}\sin^m x\,dx,
\]
where $\sin^0x$ means the constant function $1$.
The integrands are continuous and are positive in the interior
when $m\ge1$, so all $I_m$ are positive. Direct integration gives
\[
 I_0=\frac\pi2,\qquad I_1=1.
\]

\step{Optional verification, Step 1: derive the integral recurrence.}
For $m\ge2$, write
$I_m=\int_0^{\pi/2}\sin^{m-1}x\sin x\,dx$.
In integration by parts take
\[
 u=\sin^{m-1}x,\quad dv=\sin x\,dx,\quad
 du=(m-1)\sin^{m-2}x\cos x\,dx,\quad v=-\cos x.
\]
The boundary term
$[-\sin^{m-1}x\cos x]_0^{\pi/2}$ is zero:
$\cos(\pi/2)=0$, and at $0$ the factor $\sin^{m-1}0$ is zero
because $m-1\ge1$. Consequently
\begin{align*}
 I_m
  &=(m-1)\int_0^{\pi/2}\sin^{m-2}x\cos^2x\,dx\\
  &=(m-1)\int_0^{\pi/2}
          \sin^{m-2}x(1-\sin^2x)\,dx\\
  &=(m-1)(I_{m-2}-I_m).
\end{align*}
Adding $(m-1)I_m$ to both sides and dividing by $m>0$ gives
\[
 I_m=\frac{m-1}{m}I_{m-2}.
\]

\step{Optional verification, Step 2: express the even and odd integrals.}
Repeated use of the recurrence, ending at $I_0$ or $I_1$, gives
for every integer $N\ge1$,
\[
 I_{2N}=\frac\pi2\prod_{k=1}^{N}\frac{2k-1}{2k},
 \qquad
 I_{2N+1}=\prod_{k=1}^{N}\frac{2k}{2k+1}.
\]
For example, the even recurrence successively contributes the factors
$(2N-1)/(2N)$, $(2N-3)/(2N-2)$, down to $1/2$.
The odd recurrence contributes $(2N)/(2N+1)$ down to $2/3$.
These are exactly the displayed finite products.

Unrolling the recurrence for $t_n$ gives
\[
 t_{N+1}
  =\prod_{k=1}^{N}\left(1-\frac1{4k^2}\right)
  =\prod_{k=1}^{N}\frac{(2k-1)(2k+1)}{(2k)^2}.
\]
The initial empty product for $t_1$ is $1$, and each new recurrence
step multiplies by the next factor, which justifies this finite-product
identity. Dividing the formulas for the positive integrals gives
\[
 \frac{I_{2N}}{I_{2N+1}}
   =\frac\pi2
     \prod_{k=1}^{N}\frac{(2k-1)(2k+1)}{(2k)^2}
   =\frac\pi2\,t_{N+1}.
\]

\step{Optional verification, Step 3: bound the integral ratio.}
On $[0,\pi/2]$, the quantity $\sin x$ lies in $[0,1]$.
Multiplication by a number in $[0,1]$ cannot increase a nonnegative
power, so
\[
 \sin^{2N+1}x\le\sin^{2N}x\le\sin^{2N-1}x.
\]
Integrating preserves these inequalities:
\[
 I_{2N+1}\le I_{2N}\le I_{2N-1}.
\]
Divide by $I_{2N+1}>0$. The integral recurrence with $m=2N+1$
gives $I_{2N+1}=(2N)/(2N+1)\,I_{2N-1}$, so
\[
 1\le\frac{I_{2N}}{I_{2N+1}}
   \le\frac{I_{2N-1}}{I_{2N+1}}
   =\frac{2N+1}{2N}=1+\frac1{2N}.
\]
Substituting $I_{2N}/I_{2N+1}=(\pi/2)t_{N+1}$ and multiplying
by $2/\pi>0$ yields
\[
 0\le t_{N+1}-\frac2\pi\le\frac1{\pi N}.
\]
For any $\varepsilon>0$, a positive integer index
$n>1+1/(\pi\varepsilon)$ has $N=n-1>1/(\pi\varepsilon)$.
Therefore
\[
 \abs{t_n-\frac2\pi}\le\frac1{\pi(n-1)}<\varepsilon.
\]
This proves the conjectured limit by Definition~7.1.

\answer{Part (a) proves existence using Section~10 alone.
For part (b), the computed terms suggest $2/\pi$.
The explicitly marked additional calculus argument verifies that
\[
 \boxed{\displaystyle\lim_{n\to\infty}t_n=\frac2\pi.}
\]
No product formula was taken on trust; the finite-product identity
and its limiting estimate were derived.}

\exercise{10.12}{A product recurrence with an exact formula}

\goal For
\[
 t_1=1,\qquad
 t_{n+1}=\left(1-\frac1{(n+1)^2}\right)t_n\quad(n\in\N),
\]
(a) establish existence of the limit, (b) guess its value,
(c) prove $t_n=(n+1)/(2n)$ by induction, and (d) use the formula
to verify the limit.

\step{Part (a), Step 1: prove positivity and monotonicity.}
For every $n\ge1$, we have $n+1\ge2$, so
\[
 0<\frac1{(n+1)^2}\le\frac14,\qquad
 \frac34\le1-\frac1{(n+1)^2}<1.
\]
The initial term is positive. If $t_n>0$, multiplying it by this
positive coefficient gives $t_{n+1}>0$. Induction proves positivity
for every term.

Because the coefficient is strictly less than $1$ and $t_n>0$,
we also have $t_{n+1}<t_n$ for every $n$. Thus the sequence is
strictly decreasing, bounded below by $0$ and above by $t_1=1$.

\step{Part (a), Step 2: apply bounded monotone convergence.}
Theorem~10.2 now gives a finite limit $L$.
Exercise~8.9 implies $0\le L\le1$.

\answer{The limit exists as a real number. Its value has not been
assumed in establishing existence.}

\step{Part (b): calculate early terms to formulate a guess.}
Successive substitutions give
\[
 t_2=\left(1-\frac14\right)1=\frac34,
\]
\[
 t_3=\left(1-\frac19\right)\frac34
     =\frac89\cdot\frac34=\frac23,
\]
\[
 t_4=\left(1-\frac1{16}\right)\frac23
     =\frac{15}{16}\cdot\frac23=\frac58,
\]
and
\[
 t_5=\left(1-\frac1{25}\right)\frac58
     =\frac{24}{25}\cdot\frac58=\frac35.
\]
The values
\[
 1,\quad \frac34,\quad \frac23,\quad \frac58,\quad \frac35
\]
fit the proposed pattern $t_n=(n+1)/(2n)$.
Since this expression equals $1/2+1/(2n)$, it suggests the limit $1/2$.
At this stage, the pattern and the limit are guesses to be proved.

\step{Part (c), Step 1: state the induction proposition and verify its base.}
Let $P_n$ be the assertion
\[
 t_n=\frac{n+1}{2n}.
\]
For $n=1$, the right side is $(1+1)/(2\cdot1)=1$, which equals $t_1$.
Thus $P_1$ holds.

\step{Part (c), Step 2: substitute the induction hypothesis.}
Assume $P_n$ for an arbitrary positive integer $n$.
The recurrence gives
\begin{align*}
 t_{n+1}
   &=\left(1-\frac1{(n+1)^2}\right)\frac{n+1}{2n}\\
   &=\frac{(n+1)^2-1}{(n+1)^2}\frac{n+1}{2n}\\
   &=\frac{n^2+2n}{(n+1)^2}\frac{n+1}{2n}\\
   &=\frac{n(n+2)}{(n+1)^2}\frac{n+1}{2n}\\
   &=\frac{n+2}{2(n+1)}\\
   &=\frac{(n+1)+1}{2(n+1)}.
\end{align*}
The cancellation in the penultimate line is legitimate because
$n>0$ and $n+1>0$, so neither cancelled factor is zero.
The final expression is precisely the formula in $P_{n+1}$.

\answer{Mathematical induction proves
\[
 \boxed{t_n=\frac{n+1}{2n}\quad\text{for every }n\in\N.}
\]}

\step{Part (d), Step 1: separate the constant term from the error.}
The proved formula gives the identity
\[
 t_n=\frac{n+1}{2n}
     =\frac n{2n}+\frac1{2n}
     =\frac12+\frac1{2n}.
\]
Since $1/n\to0$, the limit laws from Section~9 give
$t_n\to1/2$.

\step{Part (d), Step 2: verify the conclusion directly with a threshold.}
For any $\varepsilon>0$, choose $N=1/(2\varepsilon)$.
Every positive integer $n>N$ satisfies $1/(2n)<\varepsilon$, and
\[
 \abs{t_n-\frac12}
 =\abs{\frac1{2n}}
 =\frac1{2n}
 <\varepsilon.
\]
This is the definition of convergence to $1/2$.

\answer{The guess in part (b) is now proved:
\[
 \boxed{\displaystyle\lim_{n\to\infty}t_n=\frac12.}
\]
The exact formula also gives the error $t_n-1/2=1/(2n)$ for every
positive integer index.}
